\section{相关工作}
\label{sec:related}
\paragraph{代码智能体及其接口。}
SWE-bench 使用真实问题与补丁配对以及可执行测试，评估仓库级问题修复 \cite{jimenez2024swebench}。SWE-agent 表明，搜索、读取、编辑和执行接口会影响智能体行为，其中也包括通过反馈拒绝部分无效编辑 \cite{yang2024sweagent}。OpenHands 提供可复用的智能体实现、行动与观察事件、沙箱执行和任务委派 \cite{wang2025openhands}。Agentless 则将定位、修复和验证组织成结构化流水线 \cite{xia2025agentless}。这些系统说明，有必要独立于模型规划策略研究基础设施。\wcode{} 关注操作之间交换的证据是否适用，并通过无模型评估将运行时行为与问题修复性能分开考察。

\paragraph{架构一致性与来源记录。}
软件反射模型（Software Reflexion Models）通过显式映射，将高层结构与从源码导出的模型连接起来，揭示两者的一致与偏离 \cite{murphy1995reflexion}。\wcode{} 将这一已有区别应用于实际运行的工作区，其中声明、语法、语义和运行时观察分别保留各自的含义。W3C PROV-DM 对实体、活动、负有责任的主体以及派生关系建模 \cite{w3c2013provdm}；PROV-AGENT 将来源采集扩展到智能体交互和工作流任务，包括基于 MCP 的采集 \cite{souza2025provagent}。本文关注这些记录何时能够支持当前软件操作。某项记录即使不满足修订、范围或权威性条件，仍可作为具有解释作用的历史记录保留。这是一种领域特定的准入契约，而不是新的来源记录标准。

\paragraph{复用与增量计算。}
Build Systems \`a la Carte 区分调度与重新构建，并通过依赖信息和验证轨迹刻画复用 \cite{mokhov2018buildsystems}。缓存结果必须满足相应输入条件才能使用，这一原则已经成熟。\wcode{} 将同样的推理应用于验证与验收，保留已映射检查、跳过检查、复用结果和新执行之间的差别。其贡献在于这些机制如何在共享的智能体接口运行时中与权威性和生产者范围相互作用，而不是新的增量构建算法。

\paragraph{仓库检索。}
RepoCoder 将仓库检索与迭代生成结合，Repoformer 则学习何时检索可能改善补全 \cite{zhang2023repocoder,wu2024repoformer}。两者都使检索质量与开销成为代码生成的重要因素。本文的交付契约提出一个互补问题：有界响应是否包含标注的身份、原始源码、当前文件元数据和符合条件的编辑输入。这些谓词能够在模型使用响应之前诊断工具层。要比较下游有效性，还需要匹配任务、预算和模型配置。

\paragraph{执行控制与对抗评估。}
CaMeL 将控制流与不可信数据分离，并使用权能（capability）机制对经工具传递的数据流执行策略 \cite{debenedetti2025camel}。AgentDojo 在智能体使用不可信工具输出时，区分合法任务效用与攻击者成功率 \cite{debenedetti2024agentdojo}。\wcode{} 的工作区检查、命令策略和受保护修改，针对的是具体的局部准入边界。选定反例有助于考察这些边界；而针对自适应提示注入的抵抗能力和信息流安全，则需要上述研究所考察的更广泛机制与结果指标。

这些研究方向中的基础机制均已有先行工作。本文研究的系统将它们共享的适用条件显式化，并贯穿源码观察、修改、验证和协调核对。评估在具体接口边界上测量这种组合。
